% 09_which.tex - which option to use, and the finished state
\section{Which option to use}
\begin{itemize}
\item \cmd{CONCEPTUAL}: when the water table stays below the soil, or for a first model.
\item \cmd{PHYSICAL}: where the water table reaches the soil, as on valley floors, wetlands and the foot of slopes. It
  takes more computing time per day.
\item \cmd{:SoilZoneVadose}: where the water table lies metres below the soil and the timing of the recharge matters.
  \cmd{RICHARDS} is the reference method. \cmd{KINEMATIC} leaves out capillary forces and always converges. \cmd{LAG} is a
  cascade of linear stores with a travel time per metre.
\end{itemize}
\cmd{\$WS/model/finished} holds the \cmd{.rvp} with all the parameters and the \cmd{.rvg} with \cmd{RICHARDS} and three
depth classes. The manual of the coupling explains these options in plain terms in its Volume~1 (User Manual, chapter
``The exchanges in plain terms''), and gives the equations in its Volume~2 (Reference, section ``The soil zone and the
water table''); the user guide (section~4.3c) lists every \cmd{:SoilZone} command. Both are in the
RavenHydroFramework-extra repository (see \cmd{bin/README.txt}).
